\scanpage{51}
\mspage{57}

\unsure{La} question, c'est si pour tout $y'\in Y'$, posant $y=p(y')$, \deleted{\unsure{et}} \deleted{\unsure{cons}}considérant le foncteur induit par $q$

\diagram{\begin{tikzcd}[ampersand replacement=\&]
X_{/y} \& X'_{/y'} \arrow[l]
\end{tikzcd}}

le morphisme can.

\[
\Gamma_{\unsure{\mathbb D|\underline D}}(X_{/y},F)\unsure{\relbar\joinrel\relbar|\longrightarrow}\Gamma_{\mathbb D}(X'_{/y'},q^*F)
\]

est iso. Dans le cas $f$ propre, on considère

\diagram{\begin{tikzcd}[ampersand replacement=\&]
X_y \arrow[d,"{\sim\ \text{costricte}}"'] \& X_{y'} \arrow[l,"\sim"'] \arrow[d,"{\sim\ \text{costricte}}"] \\
X_{/y} \& X'_{/y'} \arrow[l]
\end{tikzcd}}

d'où la conclusion puisque par \unsure{Dèr|Dér} 5 \unsure{g|G}, dans

\diagram{\begin{tikzcd}[ampersand replacement=\&]
\Gamma(X_y) \arrow[r] \& \Gamma(X_{y'}) \\
\Gamma(X_{/y}) \arrow[u] \arrow[r] \& \Gamma(X'_{/y'}) \arrow[u]
\end{tikzcd}}

les flèches verticales et l'hor. sup. sont iso -- donc aussi l'hor. inf. \deleted{\unsure{Dans}}

Dans le cas $p$ lisse, considérons du côté $Y/Y'$ le carré de chang$^{\mathrm{ts}}$ de base

\diagram{\begin{tikzcd}[ampersand replacement=\&]
Y \& Y' \arrow[l] \\
Y_{/y} \arrow[u] \& Y'_{/y'} \arrow[u] \arrow[l]
\end{tikzcd}}
% archive sheet number: 50
